% TOP_PROBLEM: {"problemId": "problem.babais-diameter-conjecture-for-finite-simple-groups", "canonicalTitle": "Babai's diameter conjecture for nonabelian finite simple groups", "releaseRank": 223, "primaryDomain": "algebra_representation_category", "primaryDomainLabel": "Algebra, representation & category theory", "displayStatus": "open_with_solved_subcases", "releaseStatus": "open_with_solved_subcases"}
% !TeX program = lualatex
% Definition and sources only; this file is not an LLM solution attempt.
\documentclass[11pt]{article}
\usepackage[margin=25mm]{geometry}
\usepackage{fontspec}
\setmainfont{Latin Modern Roman}
\usepackage{amsmath,amssymb,mathtools}
\usepackage{unicode-math}
\setmathfont{Latin Modern Math}
\usepackage{xurl,hyperref}
\hypersetup{colorlinks=true,urlcolor=blue}
\setcounter{secnumdepth}{0}
\setlength{\parindent}{0pt}
\setlength{\parskip}{5pt}
\setlength{\emergencystretch}{3em}
\begin{document}
\section*{\texorpdfstring{Babai's diameter conjecture for nonabelian finite simple groups}{Babai's diameter conjecture for nonabelian finite simple groups}}\label{223}
\subsection{Definitions and mathematical statement}
A nonabelian finite simple group has no nontrivial proper normal subgroup and is not commutative. For a generating set $X$, its undirected Cayley graph joins $g$ to $gx$ for $x\in X\cup X^{-1}$. Its diameter is the largest, over $g\in G$, of the shortest word length representing $g$ in these steps. The conjecture asks for one absolute $c>0$ such that
\[
 \operatorname{diam}\operatorname{Cay}(G,X\cup X^{-1})\le(\ln|G|)^c
\]
for every such $G$ and every generating set. Constants and the exponent may not depend on the group family, rank, or generators. Results requiring a special element in $X$ are subcases rather than the full assertion.
\par\smallskip\textit{Formulation and concept references:} \href{https://arxiv.org/pdf/2203.03323v1}{[S1]}.

\subsection{Short English statement}
Can every element of every nonabelian finite simple group be expressed in any generating set using only a polylogarithmic number of generators and inverses?
\subsection{Sources}
\begin{itemize}
\item[S1]Martino Garonzi, Zoltán Halasi and Gábor Somlai, On the diameter of Cayley graphs of classical groups with generating sets containing a transvection; arXiv2203.03323v1 March7,2022.
\url{https://arxiv.org/pdf/2203.03323v1}.
\textit{Formulation source.}
\end{itemize}
\end{document}
